\scanpage{12}
\mspage{8}
% source: top, cancelled first lemma
\deleted{Proposition}

\deleted{\uline{Lemme} a) Soit $a\in\unsure{L|\mathscr L}$ un élément \unsure{non terminal|non minimal}, en sorte que $I_{<a}$ est défini. Alors}
\[
\deleted{\mathrm{cos}\!\!\unsure{-}\!{\unsure{w|W|\psi}}^{\unsure{\circ|0}}(I_{<a})=\{a\}\cup\operatorname{Omb}^{\circ}(I_{>a})}
\]

% source: upper-middle, new lemma
\uline{Lemme} $\unsure{M|\mathscr M}$ est très dense dans $\widehat M$, i.e. pour $X,Y\in\widehat M$, pour qu'on ait $X\unsure{\mid\circ\mid|\mid\bullet\mid}Y$, il faut que pour tout $X'\overset{\circ}{\ll}X$, $Y'\overset{\circ}{\ll}Y$, avec $X',Y'\in M$, on ait $X'\mid\circ\mid Y'$.
% source: left margin adjacent to lemma and proof
% margin note written sideways; slanted
\margin{\uline{NB} \unsure{Il suffit} de \unsure{le|la} vérifier en prenant $X'\overset{\circ}{\ll}X$, et \unsure{pareil} $Y$.}

% source: middle, first proof paragraph
\deleted{Supp.} Il n'y a de question \deleted{que \unsure{si jamais}}\inserted{si $X$ est un} \deleted{\unsure{au moins un} des deux élém $X,Y$ est un} $I_{<a}$ ou un $I_{>b}$. Supposons p.ex.
\[
X=I_{<a},
\]
\deleted{et supposons d'abord $Y\in M$,} Alors \deleted{\unsure{X est}} $X\mid\circ\mid Y$ sss $I_{<a}<Y$ i.e. $a\leq Y$, \deleted{et on a pour $X\overset{\circ}{\ll}I_{<a}$,}\inserted{montrons qu'il suffit} $\forall X'\overset{\circ}{\ll}I_{<a}$, on ait $X'\mid\circ\mid Y$.

\deleted{\unsure{c.à.d.} $X\unsure{\leq|<}a$, \unsure{alors}}

% source: lower-middle, choice of X prime
Prenons d'abord $X'=\deleted{[a',a]}I_{\inserted{a',a}}$, pour $a'<a$ (tel $a'$ existe, car $a$ non minimal). On aura donc $Y\mid\circ\mid I_{a',a}$ donc $Y\geq a$ \uline{ou} $Y\leq a'$. Mais si on avait $Y\leq a'<a$, on aurait $Y\overset{\circ}{\ll}X=I_{<a}$, et prenant $X'=Y$, on trouve $Y\mid\circ\mid Y$, absurde.

% source: bottom, cancelled continuation
\deleted{\unsure{Si} On prouve de \unsure{même}, pour $X'$ fixé, que $X'\mid\circ\mid Y$ sss $\unsure{Y}$}

% source: bottom, corollary
\uline{Cor.} L'ens\inserted{(Collection ou cl.)} des supports $\Sigma\widehat M$ est can. isom à l'ens des supports $\Sigma M$.
% archive sheet number: 11
